\documentclass[11pt,a4paper]{article}
% Common preamble for the hint sheets (\input from each file)
\usepackage[utf8]{inputenc}
\usepackage[T1]{fontenc}
\usepackage{lmodern}
\usepackage{amsmath,amssymb,amsthm,mathtools}
\usepackage[a4paper,margin=2.3cm]{geometry}
\usepackage{enumitem}
\usepackage{xcolor}
\usepackage[hidelinks]{hyperref}
\renewcommand{\leq}{\leqslant}
\renewcommand{\le}{\leqslant}
\renewcommand{\geq}{\geqslant}
\renewcommand{\ge}{\geqslant}
\setlength{\parindent}{0pt}
\setlength{\parskip}{0.35em}
% \hintstitle{Sheet title}{Class work / Homework}
\newcommand{\hintstitle}[2]{\begin{center}{\LARGE\bfseries Hints: #1\par}\vspace{0.5em}{\large #2\par}\end{center}
\noindent\textit{The hints to each problem are cumulative. Read one, try again, and open the next only if you are still stuck.}\par\medskip\hrule\medskip}
\newcommand{\problem}[1]{\par\bigskip{\large\bfseries Problem #1}\par\nopagebreak\smallskip}
\newcommand{\subpart}[1]{\par\smallskip\textbf{(#1)}\par\nopagebreak}
\newcommand{\hint}[1]{\par\smallskip\noindent\textcolor{black!70}{\textbf{Hint #1.}}\enspace}
\newcommand{\note}[1]{\par\smallskip\noindent\textit{Note.} #1\par}
\providecommand{\lt}{<}
\providecommand{\gt}{>}
\providecommand{\tightlist}{\setlength{\itemsep}{0pt}\setlength{\parskip}{0pt}}
\emergencystretch=3em

\begin{document}
\hintstitle{Euclidean geometry -- 1. Vectors, angles, simplices and volumes in $\mathbb R^n$}{Homework}

\problem{1}
\hint{1} Realise the simplex as $A_i=\frac a{\sqrt2}e_i$ ($i=0,\dots,n$) in the hyperplane $x_0+\dots+x_n=\frac a{\sqrt2}$ of $\mathbb R^{n+1}$: all edges equal $a$, and the centre is $G=\frac a{\sqrt2}\cdot\frac1{n+1}\mathbf1$.
\hint{2} $|A_i-G|^2=\frac{a^2}2\bigl(1-\frac1{n+1}\bigr)$ gives $R$, and $\langle A_i-G,A_j-G\rangle=-\frac{a^2}{2(n+1)}$ gives the cosine $-\frac1n$. The centroid divides each median in the ratio $n:1$, so the distance from $G$ to a facet is $r=\frac Rn$.
\hint{3} The Gram matrix of $A_i-A_0$ ($i\ge1$) is $\frac{a^2}2(E+J)$ with determinant $\bigl(\frac{a^2}2\bigr)^n(n+1)$; Theorem 1 gives the volume. The outer normal of the facet opposite $A_i$ is parallel to $G-A_i$, so the outer normals of two facets make the angle $\arccos(-\frac1n)$ and the dihedral angle is $\pi-\arccos(-\frac1n)=\arccos\frac1n$.

\problem{2}
\hint{1} Expand $|XA_i|^2=|(X-G)+(G-A_i)|^2$ and use $\sum_i(G-A_i)=0$.
\hint{2} Put $X=O$: $(n+1)R^2=(n+1)|OG|^2+\sum_i|GA_i|^2$. Put $X=A_j$ and sum over $j$: $2\sum_{i<j}|A_iA_j|^2=2(n+1)\sum_i|GA_i|^2$.
\hint{3} Combine the two identities. Since $|OG|^2\ge0$, $\sum_{i<j}|A_iA_j|^2\le(n+1)^2R^2$ with equality iff $O=G$; for a tetrahedron in the unit sphere the bound is $16$.

\problem{3}
\hint{1} Take the vertices of the icosahedron $(0,\pm1,\pm\varphi)$ and their cyclic permutations, $\varphi=\frac{1+\sqrt5}2$. The six diagonals are spanned by $(0,1,\pm\varphi)$, $(1,\pm\varphi,0)$, $(\pm\varphi,0,1)$; all have $|v|^2=1+\varphi^2=\varphi+2$, and every scalar product of two of them is $\pm\varphi$ or $\pm(1-\varphi^2)=\mp\varphi$. Since $\frac{\varphi}{\varphi+2}=\frac1{\sqrt5}$, the lines are equiangular.
\hint{2} The $28$ vectors have coordinate sum $0$, so they lie in the hyperplane $\mathbf1^\perp\cong\mathbb R^7$, and $|v|^2=24$. Two of them with disjoint pairs of $3$'s have scalar product $-6-6+4=-8$; with one common position, $9-3-3+5=8$. So all angles have cosine $\pm\frac13$.
\hint{3} With $n=2$ (three lines at $60^\circ$), $n=3$ (six lines) and $n=7$ ($28$ lines) the bound $\frac{n(n+1)}2$ is attained.

\problem{4}
\hint{1} Write the ellipsoid as $x^TQx=1$ with $Q$ positive definite, and let $u_1,\dots,u_n$ be the orthonormal directions of $OA_i$, so $A_i=t_iu_i$ with $t_i^{-2}=u_i^TQu_i$.
\hint{2} In the orthonormal basis $u_i$ the hyperplane through $A_1,\dots,A_n$ is $\sum_i\frac{x_i}{t_i}=1$, so its distance from $O$ is $\Bigl(\sum_i t_i^{-2}\Bigr)^{-1/2}$.
\hint{3} $\sum_iu_i^TQu_i=\operatorname{tr}Q$ for every orthonormal basis, so the distance is $(\operatorname{tr}Q)^{-1/2}$, independent of the choice.

\problem{5}
\hint{1} Answer: $N=3m^2$. Example: $(m,m,m)$, $(m,-m,-m)$, $(-m,m,-m)$, $(-m,-m,m)$.
\hint{2} The centre of the tetrahedron is the centre of the sphere, so the vertices $v_1,\dots,v_4$ satisfy $\sum v_i=0$, $|v_i|^2=N$ and $\langle v_i,v_j\rangle=-\frac N3$. Look at the integer vectors $w_j=v_1+v_j$, $j=2,3,4$.
\hint{3} They are pairwise orthogonal and $|w_j|^2=\frac{4N}3=L^2$. The integer matrix $W$ with rows $w_j$ has $WW^T=L^2E$, so $\det W=\pm L^3$ is an integer; together with $L^2\in\mathbb Z$ this gives $L\in\mathbb Z$. Finally $N=\frac34L^2$ is an integer only for even $L$, so $N=3\bigl(\frac L2\bigr)^2$.

\problem{6}
\hint{1} Answer: exactly the odd $n$. For odd $n$ take $J-E$: rows have $n-1$ ones (even self-products), two rows share $n-2$ ones (odd products).
\hint{2} For even $n$ reduce modulo $2$: $MM^T\equiv J-E$, and $\det(J-E)=(-1)^{n-1}(n-1)$ is odd, so $M$ is invertible over $\mathbb F_2$.
\hint{3} But over $\mathbb F_2$ we have $r\cdot r\equiv r\cdot\mathbf1$ for every row $r$, so every row is orthogonal to $\mathbf1$, i.e.\ $M\mathbf1\equiv0$: a contradiction.

\problem{7}
\hint{1} Four points of a plane form a degenerate $3$-simplex, so the Cayley--Menger determinant of Theorem 1 (the bordered $5\times5$ matrix with the squared distances) vanishes.
\hint{2} The square of an odd integer is $\equiv1\pmod8$. So modulo $8$ the bordered matrix is $J-E$ of size $5$.
\hint{3} $\det(J-E)=(-1)^4\cdot4=4\not\equiv0\pmod8$, a contradiction.

\problem{8}
\hint{1} Let $P_i=v_iv_i^T$. Then $(v_i\cdot v_j)^2=\operatorname{tr}(P_iP_j)=\langle P_i,P_j\rangle$ (the scalar product of matrices), $\langle P_i,E\rangle=1$ and $\langle E,E\rangle=n$.
\hint{2} Hence $n(v_i\cdot v_j)^2-1=n\bigl\langle P_i-\frac En,P_j-\frac En\bigr\rangle$: a Gram matrix, so positive semidefinite.
\hint{3} The sum of all entries of a positive semidefinite matrix is $\mathbf1^TG\mathbf1\ge0$: $m(n-1)+m(m-1)(n\alpha^2-1)\ge0$, which is the bound. For the examples: $n=2,\alpha=\frac12$ gives $3$; $n=3,\alpha=\frac1{\sqrt5}$ gives $6$; $n=7,\alpha=\frac13$ gives $28$.

\problem{9}
\hint{1} Look for equilateral triangles with side $2\sqrt2$: three points that pairwise differ in exactly two coordinates. For instance, all neighbours of one point $x$ of the cube (points differing from $x$ in exactly one coordinate) pairwise differ in exactly two coordinates.
\hint{2} For $x\in\{\pm1\}^n$ let $N(x)$ be its $n$ neighbours. If $|B\cap N(x)|\ge3$ for some $x$ (not necessarily in $B$), we are done.
\hint{3} Double counting: every point of $B$ lies in $N(x)$ for exactly $n$ points $x$, so $\sum_x|B\cap N(x)|=n|B|>2^{n+1}=2\cdot2^n$, and some $x$ has $|B\cap N(x)|\ge3$.

\problem{10}
\hint{1} Answer: $d(k,n)=kn$ if $k,n\ge2$, and $d(k,n)=k+n$ if $k=1$ or $n=1$. For $k=1$, ``$1$-generic'' means ``not contained in a hyperplane'', and $n+1$ affinely independent points suffice. For $n=1$ the hyperplanes are points, a set is $k$-generic iff it has at least $k+1$ points, and $k+1$ points suffice.
\hint{2} Lower bound for $k,n\ge2$: put $k$ points on each of the $n$ coordinate axes, away from the origin. A hyperplane not through $0$ contains at most one point of each axis and misses at least $n(k-1)\ge k$ points; one through $0$ cannot contain all axes and misses the $k$ points of some axis. Removing any point $p$ of axis $i$ spoils this: the hyperplane through $0$ containing all other axes misses only $k-1$ points.
\hint{3} Upper bound: let $T\subseteq S$ be a minimal $k$-generic subset. Minimality gives, for every $p\in T$, a hyperplane $H_p\not\ni p$ that misses exactly $k$ points of $T$. Build a chain: $p_1\in T$, $L_1=H_{p_1}$; then $p_{j+1}\in T\cap L_j$ and $L_{j+1}=L_j\cap H_{p_{j+1}}$. Since $p_{j+1}\in L_j\setminus L_{j+1}$, the dimension drops at every step, so after at most $n$ steps $T\cap L$ has at most one point, and $T$ is covered by at most $n$ sets $T\setminus H_{p_j}$ of size $k$ plus at most one point: $|T|\le kn+1$.
\hint{4} Suppose $|T|=kn+1$. Then the $n$ sets $C_j=T\setminus H_{p_j}$ are disjoint, and there is a point $q\in T$ lying in all $H_{p_j}$ with $\bigcap H_{p_j}=\{q\}$. Every other point of $T$ lies in exactly one $C_j$, i.e.\ in all $H_{p_i}$, $i\ne j$, which intersect in a line through $q$. So $T$ is exactly the configuration of Hint~2 (with $q$ as origin) plus the point $q$, and $T\setminus\{q\}$ is still $k$-generic, contradicting minimality.
\end{document}
